\chapter{An algebraic proof of the fundamental theorem of algebra}
\label{app:algebraic-fta}

This proof replaces an appeal to the fundamental theorem of algebra as a prior theorem. Its analytic inputs are only the intermediate value theorem for real continuous functions and the existence of a nonnegative square root of a nonnegative real number. The ordered real field, the complex field $\C=\R[i]$ with $i^2=-1$, and coefficientwise complex conjugation are already constructed. No complex logarithm, exponential, winding number, transcendence theorem or statement that $\C$ is algebraically closed is used.

The algebraic inputs are the abstract splitting-field construction in Entry~\ref{K07.12.splitting_existence_uniqueness}, characteristic-zero separability in Entry~\ref{K07.12.characteristic_zero_separability}, the finite Galois theorem in Entries~\ref{K07.14.splitting_criterion}, \ref{K07.14.fundamental_correspondence} and~\ref{K07.14.degree_index}, and Sylow existence, including smaller prime-power subgroups, in Entry~\ref{K07.06.sylow_existence}. All these proofs apply to abstract fields and finite groups and were established independently of any complex root-existence assertion. In particular, this appendix does not call on a representation-theoretic diagonalization over $\C$ that might itself have used the theorem being proved.

\begin{lemma}\label{lem:real-no-odd-extension}
Every odd-degree real polynomial has a real zero, and $\R$ has no nontrivial finite extension of odd degree.
\end{lemma}
\begin{proof}
For an odd-degree polynomial $f(X)=a_dX^d+\cdots+a_0$ with $a_d\ne0$, divide by $a_d$. If $M=1+\sum_{j<d}|a_j/a_d|$, then for $|x|>M$,
\[\left|\sum_{j<d}(a_j/a_d)x^j\right|
 \leq\left(\sum_{j<d}|a_j/a_d|\right)|x|^{d-1}<|x|^d.
\]
Thus the monic polynomial is positive for $x>M$ and negative for $x<-M$, because $d$ is odd. Polynomials are continuous, so the intermediate value theorem gives a zero.

Suppose $E/\R$ were finite of odd degree and contained $a\notin\R$. By the tower law, the degree of the irreducible minimal polynomial of $a$ divides $[E:\R]$ and is therefore odd. It is greater than one, yet the first paragraph gives it a real zero. The factor theorem would then make this irreducible polynomial linear, a contradiction. Hence every element of $E$ belongs to $\R$ and $E=\R$.
\end{proof}

\begin{lemma}\label{lem:complex-square-root-algebraic}
Every complex number has a square root in $\C$. Consequently every quadratic polynomial over $\C$ splits, and $\C$ has no quadratic field extension.
\end{lemma}
\begin{proof}
Write $z=a+bi$. If $b=0$ and $a\geq0$, take its nonnegative real square root; if $b=0$ and $a<0$, use $i\sqrt{-a}$. These cases include $z=0$. If $b\ne0$, let $r=\sqrt{a^2+b^2}$. Then $r>|a|$, so
\[u=\sqrt{(r+a)/2}>0,\qquad v=\frac{b}{2u}\]
are defined real numbers. Since $(r-a)(r+a)=b^2$, one has $v^2=(r-a)/2$. Therefore $u^2-v^2=a$ and $2uv=b$, proving $(u+vi)^2=z$.

For $AX^2+BX+C$ with $A\ne0$, the quadratic formula is valid because $2A\ne0$ in characteristic zero and $B^2-4AC$ has the square root just constructed. It gives both roots, including a repeated one when the discriminant is zero. If $E/\C$ had degree two, any $a\in E\setminus\C$ would have a degree-two minimal polynomial by the tower law. Such a polynomial has a root in $\C$ and is reducible, a contradiction.
\end{proof}

\begin{theorem}[Fundamental theorem of algebra]\label{thm:algebraic-fta}
Every nonconstant polynomial in $\C[X]$ has a complex root. Equivalently, $\C$ is algebraically closed, and every polynomial of degree $n\geq1$ is a product of $n$ linear factors over $\C$, with multiplicities counted.
\end{theorem}
\begin{proof}
Let $f\in\C[X]$ have positive degree. Write $\overline f$ for coefficientwise conjugation and put
\[h(X)=(X^2+1)f(X)\overline f(X).\]
This polynomial has real coefficients: conjugation fixes $X^2+1$ and interchanges the two other factors. By the abstract splitting-field theorem there is a finite splitting field $L/\R$ for $h$. Its existence has not assumed that any of these roots are already complex numbers. The distinct irreducible factors of $h$ are separable in characteristic zero; their product is separable with the same splitting field. Consequently $L/\R$ is finite Galois.

There is a root $j\in L$ of $X^2+1$. It is not real, since the square of a real number is nonnegative. The field $T=\R(j)$ is therefore a degree-two extension of $\R$, and
\[\iota:\C\longrightarrow T,\qquad a+bi\longmapsto a+bj\]
is an $\R$-isomorphism: the defining multiplication relation is the same, and $1,j$ are real-linearly independent. We use this explicit copy of $\C$ inside $L$; we have not silently assumed the original complex field contained the abstract splitting field. In particular, by the preceding lemma $T$ has no quadratic extension.

Let $G=\Gal(L/\R)$ and choose a Sylow $2$-subgroup $P$. Its fixed field has degree $[G:P]$ over $\R$, which is odd. Lemma~\ref{lem:real-no-odd-extension} forces that fixed field to be $\R$. The Galois correspondence then gives $P=G$. Thus $G$ is a finite $2$-group.

Let $H=\Gal(L/T)$. If $H\ne1$, it has order $2^m$ with $m\geq1$. By the prime-power subgroup existence theorem it has a subgroup $J$ of order $2^{m-1}$. The fixed field $L^J$ contains $T$, and the degree-index formula gives
\[[L^J:T]=[H:J]=2.\]
This contradicts the absence of quadratic extensions of $T$. Hence $H=1$ and $L=T$.

The polynomial $h$ splits in $T$. Under $\iota$, the transported polynomial $\iota(f)$ divides $h$ in $T[X]$. By unique factorization in the field polynomial ring, every positive-degree divisor of a product of linear polynomials has a linear factor; therefore $\iota(f)$ has a zero in $T$. Applying $\iota^{-1}$ gives a zero of $f$ in $\C$. Repeated division by the corresponding linear factor lowers the degree by one and proves the complete linear factorization statement. Conversely such factorizations plainly give the root property, proving the equivalences.
\end{proof}

\begin{remark}
The proof is the Sylow-$2$ argument in Dummit--Foote, Section~14.6, Theorem~35, second proof (printed~617, PDF~631), with its real/quadratic preliminary facts on printed~615--616. Stewart, Section~23.3, Theorem~23.12 and Corollary~23.13, gives the same route in a more general ordered-field setting. Artin, Section~15.10 (printed~471, PDF~483), instead outlines a winding-number argument; that outline is not being substituted for a proof here. This appendix is independent of the later special Lindemann proof of $\pi$ transcendence, so its use as a foundation for complex splitting fields creates no transcendence cycle.
\end{remark}
